% `tex.web` §307 and §314: fifteen kinds of token list, each named in an error
% context by the tag §314 gives it.
%
%     output_text: print_nl("<output> ");
%     every_par_text: print_nl("<everypar> ");
%     every_math_text: print_nl("<everymath> ");
%     …
%
% Ours knew three of them and showed the rest as `<to be read again> `, which
% is not merely a wrong label: the tag is how a reader tells "this came from
% your \everyvbox" from "this was put back", and §317 measures the tag's width
% when it decides where to cut the line under it — so a mis-named level also
% cuts its own text in the wrong place.
%
% The other half of this case is §311's elision, which is why every error here
% is deliberately two levels deep:
%
%     if (base_ptr=input_ptr) or bottom_line or (nn<error_context_lines) then
%       @<Display the current context@>
%     else if nn=error_context_lines then
%       begin print_nl("..."); incr(nn); end;
%
% and "Display the current context" then declines to print a spent backed-up
% list anyway. The order is the point: such a level prints nothing but has
% already been counted, so with `\errorcontextlines` at its default of 0 it
% becomes the `...` rather than silence. Each `\hbox{` here leaves one behind —
% `scan_spec` backs up the token that was not `to` or `spread` — and we skipped
% it outright, so every one of these contexts lost its `...`.
\catcode`\{=1 \catcode`\}=2 \catcode`\$=3 \catcode`\#=6
\everypar={\undefA}
\everyhbox={\undefB}
\everyvbox={\undefC}
\everymath={\undefD}
\everydisplay={\undefE}
\everycr={\undefF}
\hsize=100pt

% Math needs its two symbol families before §1195 will let a formula stand;
% without them the formula is deleted and `\everymath` never runs a second time.
\font\sy=cmsy10 \textfont2=\sy \scriptfont2=\sy \scriptscriptfont2=\sy
\font\ex=cmex10 \textfont3=\ex \scriptfont3=\ex \scriptscriptfont3=\ex
\font\rm=cmr10 \textfont0=\rm \scriptfont0=\rm \scriptscriptfont0=\rm \rm
\font\mi=cmmi10 \textfont1=\mi \scriptfont1=\mi \scriptscriptfont1=\mi

\setbox0\hbox{x}
\setbox1\vbox{y}
\noindent z\par
\setbox2\hbox{$a$}
% One `\cr` only: the preamble's. A row's `\cr` leaves a spent `v` template
% on TeX's stack — §325's while loop names `v_template` as the one kind it will
% not pop early — and §311 counts that level, so the context of the `\everycr`
% after a *row* carries a `...` we do not yet produce. That is a gap in how
% this engine models an alignment's input stack, not in §314's names, and it
% belongs to whatever closes that gap rather than here.
\halign{#\cr}
\message{PROBE:done}
\end
